\chapter{Frames and Erasures}
\label{ch:erasures}
\fancyhead[R]{\small\textit{Chapter 10: Frames and Erasures}}

\section{Overview: Closing the Loop on Redundancy}
\label{sec:ch10-overview}

Chapter~1 offered two promises: redundant frames are robust to noise, and
robust to erasures. Chapter~5 fulfilled the first promise --- the noise
stability analysis showed the canonical dual minimizes expected squared
reconstruction error under additive coefficient noise. This chapter
fulfills the second.

An \emph{erasure} is a coefficient that never arrives at all, not merely
one that is corrupted by noise. In communication, some transmitted symbols
are simply lost in transit; in sensor networks, some sensors fail
entirely; in data storage, some blocks become unreadable. In all of these
settings, the question is not ``how noisy was the measurement?'' but
``can I still reconstruct at all, and how accurately?''

We will see that frames answer this question cleanly, and that the
geometry developed in Chapters~8--9 (frame potential, equiangularity)
provides the right tools for understanding which frames are most
erasure-robust. The chapter closes with a precise connection back to
Chapter~9: equiangular tight frames are not only optimal for coherence
minimization (the Grassmannian packing problem), they are also optimal
for worst-case erasure robustness.

\section{The Erasure Model}
\label{sec:erasure-model}

\begin{definition}[Erasure pattern and sub-frame]
\label{def:erasure-pattern}
Given a frame $\{f_i\}_{i=1}^m$ for $V=\F^n$, an \emph{erasure pattern}
of size $k$ is a subset $E\subset\{1,\ldots,m\}$ with $|E|=k$. The
\emph{surviving index set} is $J=\{1,\ldots,m\}\setminus E$, and the
corresponding \emph{sub-frame} is $\{f_i\}_{i\in J}$.

In the erasure model, the receiver obtains only the frame coefficients
$\{c_i=\inner{x}{f_i}\}_{i\in J}$ (the $k$ coefficients indexed by $E$
are lost), and must reconstruct $x$ from these $m-k$ surviving
measurements.
\end{definition}

\begin{definition}[$k$-erasure robustness]
\label{def:k-erasure-robustness}
A frame $\{f_i\}_{i=1}^m$ for $\F^n$ is \emph{$k$-erasure robust} if,
for every erasure pattern $E$ with $|E|=k$, the sub-frame
$\{f_i\}_{i\in J}$ (where $J=\{1,\ldots,m\}\setminus E$) still spans
$\F^n$. The \emph{erasure number} of the frame is
\[
  \epsilon(F) := \max\setdef{k}{\text{the frame is }k\text{-erasure robust}},
\]
the maximum number of simultaneous erasures the frame can \emph{guarantee}
to survive --- i.e., every pattern of that many erasures leaves a spanning
sub-frame.
\end{definition}

\begin{remark*}
Spanning is the correct condition because, by Theorem~2.2.1, a finite
collection spans $\F^n$ if and only if it is a frame --- and we need the
sub-frame to be a frame so that the canonical reconstruction formula
$x = \tilde F_J (F_J^* x)$ (via the sub-frame's canonical dual) is valid.
A non-spanning sub-frame has a direction $x_0$ with
$\inner{x_0}{f_i}=0$ for all $i\in J$; no number of surviving
measurements can recover $x_0$'s component, and reconstruction is
impossible.
\end{remark*}

\begin{proposition}[Erasure number and redundancy]
\label{prop:erasure-redundancy}
For a frame $\{f_i\}_{i=1}^m$ for $\F^n$:
\begin{enumerate}[label=(\roman*)]
  \item $\epsilon(F)\le m-n$ (can tolerate at most $m-n$ erasures).
  \item $\epsilon(F)=0$ if $m=n$ (no erasure tolerance, since any single
        erasure leaves $n-1$ vectors, which cannot span $\F^n$).
  \item The erasure number depends on the \emph{geometry} of the frame,
        not merely the redundancy $m-n$: two frames with the same $m$
        and $n$ can have different erasure numbers.
\end{enumerate}
\end{proposition}

\begin{proof}
(i) After $m-n+1$ erasures only $n-1$ vectors remain, which cannot span
an $n$-dimensional space.
(ii) With $m=n$, erasing any single vector leaves $n-1$ vectors in $\F^n$,
insufficient to span.
(iii) Demonstrated in Examples~\ref{ex:square-triangle-comparison} and
\ref{ex:pentagon-erasure} below: the square and pentagon frames in $\R^2$
have the same $n=2$ but very different erasure numbers despite both being
tight.
\end{proof}

\section{Reconstruction After Erasure}
\label{sec:reconstruction-after-erasure}

If the sub-frame $\{f_i\}_{i\in J}$ spans $\F^n$, reconstruction is
always possible in principle: use the canonical dual of the sub-frame.
But the sub-frame's dual depends on $J$, so the receiver must either
recompute it for each erasure pattern --- a computation that scales with
the number of possible patterns --- or use the original dual and accept
some approximation error. For tight frames, the second option is much
more transparent than it sounds.

\begin{proposition}[Exact reconstruction after erasure]
\label{prop:exact-recon-after-erasure}
If $\{f_i\}_{i\in J}$ is a (sub-)frame for $\F^n$ with frame operator
$S_J = \sum_{i\in J} f_if_i^*$, then
\[
  x = S_J^{-1}\!\left(\sum_{i\in J}\inner{x}{f_i}f_i\right)
  = \sum_{i\in J}\inner{x}{f_i}(S_J^{-1}f_i),
\]
a reconstruction using only the surviving coefficients.
\end{proposition}

\begin{proof}
Immediate from the canonical dual formula (Proposition~3.8.1 of
Chapter~3, restated for the sub-frame $\{f_i\}_{i\in J}$). Since the
sub-frame spans $\F^n$ by the $k$-erasure robustness assumption, $S_J$
is invertible (Proposition~3.4.2(iii)).
\end{proof}

\subsection{Sub-Frame Operator Decomposition}
\label{subsec:subframe-decomposition}

The sub-frame operator $S_J$ has a clean relationship to the full frame
operator $S$ that is worth making explicit.

\begin{proposition}[Sub-frame operator]
\label{prop:subframe-operator}
Let $E=\{1,\ldots,m\}\setminus J$ be the erased indices. Then
\[
  S_J = S - S_E, \qquad\text{where } S_E = \sum_{i\in E}f_if_i^*
\]
is the ``erased contribution'' to the frame operator. In particular, if
$\{f_i\}_{i=1}^m$ is tight with bound $A$, then
\begin{equation}
\label{eq:subframe-tight}
  S_J = AI - S_E.
\end{equation}
\end{proposition}

\begin{proof}
Direct from $S = \sum_{i=1}^m f_if_i^* = \sum_{i\in J}f_if_i^* +
\sum_{i\in E}f_if_i^* = S_J + S_E$.
\end{proof}

\subsection{Error Formula for Tight Frames}
\label{subsec:tight-erasure-error}

For tight frames, a particularly clean scenario arises when the receiver
uses the \emph{original} tight reconstruction formula
$\hat x = \tfrac{1}{A}\sum_i c_i f_i$ but simply treats erased
coefficients as zero, rather than recomputing the sub-frame dual.

\begin{proposition}[Tight-frame erasure error]
\label{prop:tight-erasure-error}
Let $\{f_i\}_{i=1}^m$ be a tight frame with bound $A$ and unit-norm
vectors, and suppose the coefficients indexed by $E$ are erased. If the
receiver applies the original tight reconstruction formula with missing
coefficients set to zero,
\[
  \hat x = \frac{1}{A}\sum_{i\in J}\inner{x}{f_i}\,f_i,
\]
then the reconstruction error is
\begin{equation}
\label{eq:tight-erasure-error}
  \hat x - x = -\frac{1}{A}\sum_{i\in E}\inner{x}{f_i}\,f_i.
\end{equation}
In particular, for a single erasure $E=\{j\}$,
\[
  \hat x - x = -\frac{\inner{x}{f_j}}{A}\,f_j,
  \qquad
  \norm{\hat x - x} = \frac{|\inner{x}{f_j}|}{A}.
\]
\end{proposition}

\begin{proof}
Since $\{f_i\}$ is tight with bound $A$, the exact reconstruction is
$x = \tfrac{1}{A}\sum_{i=1}^m \inner{x}{f_i}f_i$ (Chapter~3,
equation~(3.5)). Subtracting:
\[
  \hat x - x = \frac{1}{A}\!\sum_{i\in J}\!\inner{x}{f_i}f_i
  - \frac{1}{A}\!\sum_{i=1}^m\!\inner{x}{f_i}f_i
  = -\frac{1}{A}\!\sum_{i\in E}\!\inner{x}{f_i}f_i.
\]
For a single erasure $E=\{j\}$, the norm is
$\|\hat x-x\|=\tfrac{|\inner{x}{f_j}|}{A}\,\|f_j\|=\tfrac{|\inner{x}{f_j}|}{A}$
using $\|f_j\|=1$.
\end{proof}

\begin{keyidea}
Proposition~\ref{prop:tight-erasure-error} has two immediate
consequences.
\begin{itemize}
  \item \textbf{Error is readable without knowing $x$}: The error
        $\|\hat x-x\|=|\inner{x}{f_j}|/A$ depends on the erased
        coefficient $c_j=\inner{x}{f_j}$, which the receiver already
        knows was supposed to arrive but did not. This gives a
        computable \emph{error certificate}: the receiver can compute
        an upper bound on the error without the original signal $x$.
  \item \textbf{Tight frames localize the error}: The error
        $-\tfrac{1}{A}c_j f_j$ points \emph{in the direction} of the
        erased frame vector $f_j$ --- the damage is confined to one
        direction, not spread throughout the space. Contrast this with
        a non-tight frame, where erasing one coefficient perturbs the
        reconstruction in an entangled, direction-dependent way.
\end{itemize}
\end{keyidea}

\section{Comparing Frames: The Square, Triangle, and Pentagon}
\label{sec:frame-comparison}

We now work through the three main running examples, establishing that
the erasure number depends critically on the frame geometry, not only on
the redundancy.

\begin{examplebox}[title={Example 10.1: Triangle vs.\ square vs.\ pentagon in $\R^2$}]
\label{ex:square-triangle-comparison}
\textbf{Triangle ($m=3$, $n=2$, redundancy $1$):} Any two of the three
vectors $f_1,f_2,f_3$ at $120^\circ$ are linearly independent (no pair
is parallel or antipodal), so any single erasure leaves two
spanning vectors. Thus $\epsilon(F)=1$. This matches the redundancy
$m-n=1$, the theoretical maximum.

\textbf{Square ($m=4$, $n=2$, redundancy $2$):} The four vectors
$(\pm1,0)$ and $(0,\pm1)$ have two antipodal pairs: $(f_0,f_2)$ and
$(f_1,f_3)$. Erasing one vector from each antipodal pair (e.g.\ $E=
\{0,2\}$) leaves $\{f_1,f_3\}=\{(0,1),(0,-1)\}$, which does
\emph{not} span $\R^2$ (both vectors lie on the $y$-axis). So
$\epsilon(F)=1$, even though redundancy $m-n=2$. \textbf{The square
frame wastes its extra redundancy}: it can handle any one erasure but
not any two, despite having twice the surplus over the triangle.
\end{examplebox}

\begin{examplebox}[title={Example 10.2: Pentagon and the role of equiangularity}]
\label{ex:pentagon-erasure}
The regular pentagon ($m=5$, $n=2$, vertices at $2\pi k/5$) has
redundancy $m-n=3$, and remarkably achieves $\epsilon(F)=3$: any
$3$-element erasure leaves only $2$ vectors, but since no two pentagon
vertices are parallel (the minimum angle between adjacent vertices is
$72^\circ$, and no two vertices are antipodal), any $2$-element subset
spans $\R^2$. Since $m-n=3$ is the maximum possible, the pentagon
\emph{achieves its theoretical maximum erasure number.}

The square, by contrast, achieves only $\epsilon(F)=1$ against its
theoretical maximum of $2$. The structural deficiency is its antipodal
pairs: two of its $\binom{4}{2}=6$ size-$2$ subsets fail to span.
The pentagon has no such pairs: all $\binom{5}{2}=10$ size-$2$
subsets are spanning.
\end{examplebox}

\begin{remark*}
The phenomenon illustrated here is general: a frame achieves
$\epsilon(F)=m-n$ (the theoretical maximum) if and only if every
$(m-k)$-subset spans $\F^n$ for all $k\le m-n$, which is equivalent to
every $n$-element subset of the frame being linearly independent. This
is the \emph{full spark} property.
\end{remark*}

\begin{definition}[Full spark]
\label{def:full-spark}
A frame $\{f_i\}_{i=1}^m$ for $\F^n$ has \emph{full spark} if every
$n$-element subset of $\{f_1,\ldots,f_m\}$ is linearly independent. A
full-spark frame achieves $\epsilon(F)=m-n$, the maximum possible
erasure number.
\end{definition}

\begin{remark*}
In the language of coding theory, the spark of a matrix $F$ is the
smallest number of columns that are linearly dependent; a full-spark
frame has spark $n+1$. The triangle has spark $3=n+1$; the pentagon has
spark $3=n+1$ (any two vectors are independent); but the square has
spark $2<n+1=3$, since two of its columns ($f_0$ and $f_2$) are linearly
dependent.
\end{remark*}

\section{Worst-Case Erasure Error and ETFs}
\label{sec:worst-case-etf}

For a tight frame that is full-spark, we can reconstruct after any
$k$-erasure pattern (with $k\le m-n$). But reconstruction quality varies
by pattern: some erasures cause larger errors than others.
Proposition~\ref{prop:tight-erasure-error} shows that a single erasure of
index $j$ causes error $|\inner{x}{f_j}|/A$, which depends on both $x$
and $j$. A natural worst-case measure for a single erasure is
\[
  \text{worst-case error} = \max_j\,\max_{\norm{x}=1}|\inner{x}{f_j}|/A
  = \frac{\max_j \norm{f_j}}{A} = \frac{1}{A}
\]
for unit-norm frames (using $\max_{\norm{x}=1}|\inner{x}{f_j}|=\norm{f_j}$
by Cauchy--Schwarz, and $\norm{f_j}=1$). All tight unit-norm frames give
the same single-erasure worst case.

The situation changes for \emph{multiple} erasures. For two simultaneous
erasures $E=\{j,k\}$ and a unit-norm $x$, the naive error
$\|\hat x-x\|=\|\tfrac{1}{A}(\inner{x}{f_j}f_j+\inner{x}{f_k}f_k)\|$
depends on both the angle between $f_j$ and $f_k$ and their alignment
with $x$. Maximizing this over $x$ gives
\[
  \max_{\norm{x}=1}\left\|\frac{1}{A}(\inner{x}{f_j}f_j+\inner{x}{f_k}f_k)\right\|
  = \frac{\lambda_{\max}(S_{\{j,k\}})}{A},
\]
where $\lambda_{\max}(S_{\{j,k\}})$ is the largest eigenvalue of
$S_{\{j,k\}}=f_jf_j^*+f_kf_k^*$. For unit-norm vectors,
$\lambda_{\max}(S_{\{j,k\}})=1+|\inner{f_j}{f_k}|$ (a standard fact
about the largest eigenvalue of a sum of two rank-one unit-norm
projections). So the worst-case $2$-erasure naive error is
$(1+|\inner{f_j}{f_k}|)/A$, and minimizing this over all pairs $(j,k)$
(to find the most robust frame) requires \emph{minimizing} the
off-diagonal inner products $|\inner{f_j}{f_k}|$ --- exactly what
coherence minimization and the Welch bound from Chapter~9 do.

\begin{theorem}[ETFs minimize worst-case erasure error]
\label{thm:etf-erasure-optimal}
Among all tight, unit-norm frames of $m$ vectors in $\F^n$ with $m>n$,
the equiangular tight frames (ETFs) minimize the worst-case
$2$-erasure naive reconstruction error. The minimum worst-case error
equals $(1+c_{\mathrm{Welch}})/A$ where $c_{\mathrm{Welch}} =
\sqrt{(m-n)/(n(m-1))}$ is the Welch bound coherence and $A=m/n$.
\end{theorem}

\begin{proof}
By the analysis above, the worst-case $2$-erasure naive error over all
pairs $(j,k)$ and all unit-norm signals is
$\max_{j\ne k}(1+|\inner{f_j}{f_k}|)/A$. This is minimized by
minimizing $\max_{j\ne k}|\inner{f_j}{f_k}|$, i.e.\ the coherence
$\mu$. By the Welch bound (Theorem~9.3.1), $\mu\ge c_{\mathrm{Welch}}$
for any tight unit-norm frame, with equality iff the frame is an ETF.
\end{proof}

\begin{keyidea}
\textbf{ETFs are optimal for erasure robustness, coherence minimization,
and Grassmannian packing --- simultaneously.} All three criteria reduce to
the same underlying condition: equal off-diagonal inner products achieving
the Welch bound. The equivalence established in Theorem~\ref{thm:etf-erasure-optimal}
shows that the three different motivations (information theory, geometry,
coding theory) converge to the same object.
\end{keyidea}

\section{NumPy Implementation}
\label{sec:numpy-erasures}

\begin{lstlisting}[caption={Erasure robustness tools in NumPy}]
import numpy as np
from itertools import combinations

def subframe_spans(F, erased_indices):
    """
    Return True if the sub-frame (columns NOT in erased_indices) spans R^n.
    """
    n, m = F.shape
    kept = [i for i in range(m) if i not in erased_indices]
    if len(kept) < n:
        return False
    return np.linalg.matrix_rank(F[:, kept]) == n

def erasure_number(F):
    """
    Maximum k such that every k-erasure pattern leaves a spanning sub-frame.
    """
    n, m = F.shape
    for k in range(1, m - n + 2):
        all_span = all(
            subframe_spans(F, list(E))
            for E in combinations(range(m), k)
        )
        if not all_span:
            return k - 1
    return m - n

def reconstruct_after_erasure(F, x, erased_indices):
    """
    Reconstruct x from surviving frame coefficients.
    Returns (x_hat, error_norm) using sub-frame canonical dual.
    """
    n, m = F.shape
    kept = [i for i in range(m) if i not in erased_indices]
    F_sub = F[:, kept]
    S_sub = F_sub @ F_sub.T
    F_sub_tilde = np.linalg.solve(S_sub, F_sub)
    c_sub = F_sub.T @ x
    x_hat = F_sub_tilde @ c_sub
    return x_hat, np.linalg.norm(x_hat - x)

def tight_erasure_error_formula(F, x, erased_indices, A):
    """
    Error using ORIGINAL tight dual with erased coefficients set to zero.
    Returns error vector via Proposition 10.3.3.
    """
    erased = np.array(erased_indices)
    c_erased = F[:, erased].T @ x
    error = -(1/A) * (F[:, erased] @ c_erased)
    return error

# -------------------------------------------------------------------
# Example 10.1: Triangle, square, pentagon
# -------------------------------------------------------------------
# Triangular frame (tight, A=3/2)
f1,f2,f3 = ([1.,0.],[-0.5,np.sqrt(3)/2],[-0.5,-np.sqrt(3)/2])
F_tri = np.array([f1,f2,f3]).T

# Square frame (tight, A=2)
angles4 = np.array([0,np.pi/2,np.pi,3*np.pi/2])
F_sq = np.vstack([np.cos(angles4), np.sin(angles4)])

# Pentagon frame (tight, A=5/2)
angles5 = 2*np.pi*np.arange(5)/5
F_pent = np.vstack([np.cos(angles5), np.sin(angles5)])

print("=== Erasure numbers ===")
for name, F in [("Triangle (m=3,n=2)", F_tri),
                ("Square   (m=4,n=2)", F_sq),
                ("Pentagon (m=5,n=2)", F_pent)]:
    n, m = F.shape
    eps = erasure_number(F)
    print(f"  {name}: epsilon(F)={eps}, m-n={m-n} "
          f"{'(full spark!)' if eps==m-n else '(sub-optimal)'}")

# -------------------------------------------------------------------
# Example 10.3: Tight-frame erasure error formula
# -------------------------------------------------------------------
x = np.array([3., -1.])
A_tri = 3/2
print("\n=== Tight-frame erasure error formula (triangle, erase f1) ===")
# Erase index 0 (f1)
err_vec = tight_erasure_error_formula(F_tri, x, [0], A_tri)
x_hat_naive = (1/A_tri) * (F_tri[:, 1:] @ (F_tri[:, 1:].T @ x))
err_actual = x_hat_naive - x
print(f"  Formula predicts: {err_vec.round(4)}")
print(f"  Actual error:     {err_actual.round(4)}")
print(f"  Match: {np.allclose(err_vec, err_actual)}")
print(f"  Error norm = |c1|/A = {abs(np.dot(x,F_tri[:,0])):.4f}/{A_tri} "
      f"= {abs(np.dot(x,F_tri[:,0]))/A_tri:.4f}")

# -------------------------------------------------------------------
# Example 10.4: Exact reconstruction after erasure (sub-frame dual)
# -------------------------------------------------------------------
print("\n=== Exact reconstruction (sub-frame dual) ===")
for i in range(3):
    x_hat, err = reconstruct_after_erasure(F_tri, x, [i])
    print(f"  Triangle, erase f{i+1}: error = {err:.2e}")

for i in range(4):
    x_hat, err = reconstruct_after_erasure(F_sq, x, [i])
    print(f"  Square, erase f{i+1}: error = {err:.2e}")

# Demonstrate square failure for 2 erasures (antipodal pair)
print("\n=== Square: which 2-erasure patterns fail? ===")
for pair in combinations(range(4), 2):
    spans = subframe_spans(F_sq, list(pair))
    print(f"  Erase {pair}: sub-frame spans = {spans}")
\end{lstlisting}

\section{Chapter Summary}
\label{sec:summary-ch10}

\begin{itemize}
  \item A frame is \textbf{$k$-erasure robust} if every erasure pattern
        of size $k$ leaves a spanning sub-frame (Definition~\ref{def:k-erasure-robustness}).
        The \textbf{erasure number} $\epsilon(F)$ is the maximum
        guaranteed-survivable number of erasures; it satisfies
        $\epsilon(F)\le m-n$ for any frame in $\F^n$.

  \item \textbf{Exact reconstruction} after $k$ erasures is always
        possible when the sub-frame spans, via the sub-frame's canonical
        dual $S_J^{-1}$ (Proposition~\ref{prop:exact-recon-after-erasure}).
        The sub-frame operator is $S_J = S - S_E$
        (Proposition~\ref{prop:subframe-operator}).

  \item For a \textbf{tight frame}, if the receiver uses the original
        tight formula with erased coefficients set to zero, the error
        is $\hat x - x = -\tfrac{1}{A}\sum_{i\in E}\inner{x}{f_i}f_i$
        (Proposition~\ref{prop:tight-erasure-error}) --- a clean,
        computable expression pointing exactly in the direction of the
        erased frame vectors.

  \item The \textbf{full spark} property ($\epsilon(F)=m-n$) requires
        every $n$-element subset to be linearly independent. The triangle
        and pentagon in $\R^2$ are full-spark; the square is not (its
        antipodal pairs have spark $2 < n+1$).

  \item \textbf{Theorem~\ref{thm:etf-erasure-optimal}}: Among tight
        unit-norm frames, ETFs minimize worst-case $2$-erasure error.
        The optimal error is $(1+c_{\text{Welch}})/A$, achieved exactly
        when $\mu = c_{\text{Welch}}$ (the Welch bound from Chapter~9).
        ETFs are simultaneously optimal for coherence, Grassmannian
        packing, and worst-case erasure robustness.
\end{itemize}

\section{Lab Exercises}
\label{sec:lab-ch10}

\begin{exercisebox}[title={Lab Exercise 10.1: Erasure Numbers}]
Implement \texttt{erasure\_number(F)} and compute it for:
(a)~the triangular frame, (b)~the square frame, (c)~the pentagon,
(d)~the ONB $\{e_1,e_2\}$ in $\R^2$. For each, compare to the theoretical
maximum $m-n$. Identify which are full-spark and which are not.
\end{exercisebox}

\begin{exercisebox}[title={Lab Exercise 10.2: Sub-Frame Reconstruction}]
Using the triangular frame and $x=(3,-1)^T$:
\begin{enumerate}[label=(\alph*)]
  \item For each single-erasure pattern, use \texttt{reconstruct\_after\_erasure}
        and confirm the error is zero (exact reconstruction).
  \item For each single-erasure pattern, also apply the original tight formula
        with the missing coefficient set to zero. Compute the error and verify
        it matches the formula $-\tfrac{1}{A}\inner{x}{f_j}f_j$.
  \item Which erasure is ``cheapest'' (smallest error in (b))? Explain in terms
        of which coefficient carries the least information about $x$.
\end{enumerate}
\end{exercisebox}

\begin{exercisebox}[title={Lab Exercise 10.3: The Square's Weakness}]
For the square frame, run \texttt{subframe\_spans} on all $\binom{4}{2}=6$
size-2 erasure patterns and identify exactly which two patterns cause failure.
What is special about the erased pairs? Now construct a modified
$4$-vector frame in $\R^2$ by slightly rotating two vectors so that no pair
is antipodal; verify that the modified frame is full-spark (\texttt{erasure\_number}~$=m-n=2$). Does the modified frame remain tight? (Use \texttt{is\_etf} from Chapter~9.)
\end{exercisebox}

\begin{exercisebox}[title={Lab Exercise 10.4: Worst-Case Erasure Error}]
For each of the triangle, square, and pentagon frames and a fixed random $x$:
\begin{enumerate}[label=(\alph*)]
  \item Compute the tight-formula error $\|\hat x-x\|$ for every possible
        single erasure, and report the maximum (worst case) over the $m$
        erasure patterns.
  \item Repeat for all $\binom{m}{2}$ double-erasure patterns (where applicable
        --- use only spanning ones for the square).
  \item Verify Theorem~\ref{thm:etf-erasure-optimal} numerically: among
        the three frames, which achieves the smallest worst-case
        double-erasure error, and is it the one with smallest coherence?
\end{enumerate}
\end{exercisebox}

\begin{exercisebox}[title={Lab Exercise 10.5 (Challenge): Full-Spark Random Frames}]
In theory, a ``generic'' frame (i.e.\ one drawn from a continuous
probability distribution on $(\F^n)^m$) is almost surely full-spark.
Verify this numerically: generate $100$ random unit-norm frames with
$(n,m)=(2,5)$ (draw Gaussian and normalize columns), and for each
compute \texttt{erasure\_number}. What fraction achieve $\epsilon(F)=m-n=3$?
Then try $(n,m)=(3,6)$. Explain briefly why full-spark is generic but
not universal.
\end{exercisebox}
\newpage
\section{Supplementary Exercises}
\label{sec:extra-ch10}

\subsection*{Theoretical Exercises}
\begin{theoexercises}
	\item Prove $\epsilon(F)\le m-n$ for any frame (a counting argument).
	Separately, give a full proof that a frame is full-spark if and
	only if every $n$-element subset of $\{f_1,\ldots,f_m\}$ is
	linearly independent.
	
	\item Prove that the full-spark property is preserved under \emph{any}
	invertible linear transformation $T$ (not merely orthogonal
	ones): if $\{f_i\}$ is full-spark, so is $\{Tf_i\}$. Contrast this
	with tightness (Chapter~4 Supplementary Exercise T2), which is
	only preserved under \emph{orthogonal} transformations.
	
	\item Prove that the set of non-full-spark configurations of $m$
	vectors in $\R^n$ has Lebesgue measure zero (as a subset of
	$\R^{n\times m}$), justifying the empirical finding of Lab
	Exercise~10.5 that full-spark is generic. (Hint: express
	``some $n$-subset is dependent'' as a finite union of zero sets
	of polynomials --- namely the relevant $n\times n$ minors --- and
	argue each such zero set has measure zero unless the polynomial
	is identically zero.)
	
	\item Using $S_J=S-S_E$ (Proposition~10.3.2) and Weyl's inequality for
	eigenvalues of a perturbed symmetric matrix, prove the computable
	lower bound $A_J\ge A-\lambda_{\max}(S_E)$ on the sub-frame's
	lower bound after erasing the index set $E$, without needing to
	verify spanning by brute force.
	
	\item Prove that the erasure number $\epsilon(F)$ is invariant under
	scaling all frame vectors by a common nonzero constant $c$:
	$\epsilon(cF)=\epsilon(F)$.
\end{theoexercises}

\subsection*{Computational Exercises}
\begin{compexercises}
	\item For the triangle, square, and pentagon, directly verify the
	equivalence from T1: check that ``full-spark'' (as computed by
	\texttt{erasure\_number}) agrees with ``every $n$-subset is
	linearly independent'' (checked directly via
	\texttt{np.linalg.matrix\_rank} on each $n$-subset).
	
	\item Apply a random invertible (non-orthogonal) $2\times2$ matrix to
	the pentagon and confirm the transformed configuration is still
	full-spark, illustrating T2.
	
	\item For the triangular frame, erase each single vector in turn,
	compute the exact sub-frame bound $A_J$, and compare to the Weyl
	lower bound $A-\lambda_{\max}(S_E)$ from T4. Report whether the
	bound is tight (achieved with equality) or merely a strict
	inequality in each case.
	
	\item Scale the pentagon by factors $c=0.5,2,10$ and confirm the
	erasure number is unchanged in every case, verifying T5.
\end{compexercises}